% Section 3: Implementation.
% Source paths in "% src:" comments.
%   src/..., tests/... and Cargo.toml are relative to /home/guarracino/git/impg at HEAD 0987752 (impg 0.5.0).
%   "impg <cmd> --help" refers to the saved help texts /tmp/claude-1000/-home-guarracino-git-impg-paper/ed3cfa77-35b3-454f-91de-3d0a63c6ad53/scratchpad/phase2/plan_help_<cmd>.txt.
%   plan.md is /home/guarracino/git/impg-paper/drafts/bioinformatics-2026-11/supporting/plan.md.
%   VERIFY_CORRECTIONS.md is /tmp/claude-1000/-home-guarracino-git-impg-paper/ed3cfa77-35b3-454f-91de-3d0a63c6ad53/scratchpad/phase1/VERIFY_CORRECTIONS.md.

\section{Implementation}
\label{sec:implementation}

\begin{figure*}[t]
\centering
\resizebox{\textwidth}{!}{%
\begin{tikzpicture}[
  font=\footnotesize,
  >={Stealth[length=1.8mm]},
  bar/.style={draw, fill=gray!15},
  plab/.style={font=\bfseries\normalsize, anchor=north west, inner sep=0pt},
  ptitle/.style={font=\footnotesize\bfseries, anchor=west},
  entry/.style={draw, rounded corners=1pt, font=\scriptsize, minimum height=4mm, minimum width=19mm, inner sep=1.5pt, anchor=west},
  fline/.style={font=\ttfamily\scriptsize, anchor=west, inner sep=1.5pt},
  blk/.style={draw, fill=gray!10, minimum height=5mm, font=\scriptsize, inner sep=1.5pt},
  obox/.style={draw, rounded corners=2pt, fill=white, align=left, font=\ttfamily\scriptsize, inner sep=2pt, anchor=north west},
  hop/.style={->, thick, red!70!black},
]
%% panel a: sequences and alignments
\node[plab] at (0,9.95) {a};
\node[ptitle] at (0.4,9.8) {Alignments as an implicit graph};
\fill[blue!20] (1.4,8.8) -- (3.4,8.8) -- (3.4,7.7) -- (1.4,7.7) -- cycle;
\fill[blue!20] (3.9,8.8) -- (5.9,8.8) -- (5.4,7.7) -- (3.4,7.7) -- cycle;
\fill[orange!35] (3.0,7.4) -- (7.6,7.4) -- (3.0,6.3) -- (7.6,6.3) -- cycle;
\draw[bar] (1.0,8.8) rectangle (8.2,9.1);
\node[anchor=east] at (0.95,8.95) {A};
\draw[bar] (1.0,7.4) rectangle (8.2,7.7);
\node[anchor=east] at (0.95,7.55) {B};
\draw[bar] (1.0,6.0) rectangle (8.2,6.3);
\node[anchor=east] at (0.95,6.15) {C};
\draw[decorate, decoration={brace, amplitude=2pt}] (3.4,9.15) -- (3.9,9.15) node[midway, above=1pt, font=\scriptsize] {insertion in A};
\node[font=\scriptsize] at (2.4,8.25) {$a_1$, strand $+$};
\node[font=\scriptsize, anchor=west] at (7.7,6.85) {$a_2$, strand $-$};
\node[font=\scriptsize] at (4.6,5.55) {$G(A)$ is implied by $A$ and never built};
%% panel b: index
\node[plab] at (9.4,9.95) {b};
\node[ptitle] at (9.8,9.8) {Index: one interval tree per target};
\node[anchor=east] at (10.9,9.05) {$T_A$};
\node[entry, fill=blue!20] (eA) at (11.0,9.05) {$a_1^{R}\to$ B};
\node[anchor=east] at (10.9,8.5) {$T_B$};
\node[entry, fill=blue!20] (eB1) at (11.0,8.5) {$a_1\to$ A};
\node[entry, fill=orange!35] (eB2) at (11.0,8.0) {$a_2^{R}\to$ C};
\node[anchor=east] at (10.9,7.45) {$T_C$};
\node[entry, fill=orange!35] (eC) at (11.0,7.45) {$a_2\to$ B};
\draw[fill=gray!5] (13.3,7.2) rectangle (18.3,9.4);
\node[font=\scriptsize, anchor=north west] at (13.35,9.38) {alignment file (PAF, TPA or 1ALN)};
\node[fline, fill=blue!20] (r1) at (13.45,8.65) {A 0 110 + B 0 100 50=10I50=};
\node[fline, fill=orange!35] (r2) at (13.45,7.75) {B 0 100 - C 0 100 100=};
\draw[->, gray!60!black] (eA.east) to[out=0, in=180] (r1.west);
\draw[->, gray!60!black] (eB1.east) to[out=0, in=180] (r1.west);
\draw[->, gray!60!black] (eB2.east) to[out=0, in=180] (r2.west);
\draw[->, gray!60!black] (eC.east) to[out=0, in=180] (r2.west);
\node[font=\scriptsize, anchor=west] at (9.8,6.95) {entry: other sequence, coordinates, strand, flag $R$, offset or record number (no CIGAR)};
\node[font=\scriptsize, anchor=west] at (9.8,6.55) {IMPGIDX2 file};
\node[blk, minimum width=9mm, anchor=west] (m1) at (9.8,6.1) {magic};
\node[blk, minimum width=13mm, right=0pt of m1] (m2) {map offset};
\node[blk, minimum width=10mm, right=0pt of m2] (m3) {names};
\node[blk, minimum width=9mm, fill=white, right=0pt of m3] (m4) {$T_A$};
\node[blk, minimum width=9mm, fill=white, right=0pt of m4] (m5) {$T_B$};
\node[blk, minimum width=9mm, fill=white, right=0pt of m5] (m6) {$T_C$};
\node[blk, minimum width=13mm, right=0pt of m6] (m7) {tree map};
\draw[->, red!70!black] (m7.south) to[out=-150, in=-30] node[below, font=\scriptsize] {decoded on first use} (m5.south);
%% panel c: two-hop closure
\node[plab] at (0,4.75) {c};
\node[ptitle] at (0.4,4.6) {Two-hop projection of $I$ ($m=2$)};
\draw[bar] (1.0,3.6) rectangle (8.2,3.9);
\node[anchor=east] at (0.95,3.75) {A};
\draw[bar] (1.0,2.2) rectangle (8.2,2.5);
\node[anchor=east] at (0.95,2.35) {B};
\draw[bar] (1.0,0.8) rectangle (8.2,1.1);
\node[anchor=east] at (0.95,0.95) {C};
\fill[red!60] (2.6,3.6) rectangle (4.4,3.9);
\fill[red!20] (3.2,3.6) rectangle (3.8,3.9);
\node[font=\scriptsize, anchor=south] at (3.5,3.92) {$I=$ A:40-70, depth 0};
\fill[red!60] (2.6,2.2) rectangle (3.8,2.5);
\draw[->, white, thick] (2.75,2.35) -- (3.65,2.35);
\node[font=\scriptsize, anchor=south west] at (4.2,2.52) {B:40-60, strand $+$, depth 1};
\fill[red!60] (2.6,0.8) rectangle (3.8,1.1);
\draw[->, white, thick] (3.65,0.95) -- (2.75,0.95);
\node[font=\scriptsize, anchor=south west] at (4.2,1.12) {C:40-60, strand $-$, depth 2};
\draw[hop] (3.2,3.55) -- node[right, font=\scriptsize] {through $a_1^{R}$} (3.2,2.55);
\draw[hop] (3.2,2.15) -- node[right, font=\scriptsize] {through $a_2^{R}$} (3.2,1.15);
\node[font=\scriptsize] at (4.6,0.3) {visited sets $V_A$, $V_B$, $V_C$ keep only new parts};
%% panel d: outputs
\node[plab] at (9.4,4.75) {d};
\node[ptitle] at (9.8,4.6) {Outputs};
\node[obox] at (9.8,4.25) {BED (merged)\\A 40 70 A:40-70 . +\\B 40 60 A:40-70 . +\\C 40 60 A:40-70 . -};
\node[obox] at (14.6,4.25) {FASTA\\>A:40-70\\ACGT...\\>B:40-60\\ACGT...\\>C:40-60\\ACGT...};
\node[obox] at (9.8,2.0) {PAF, one row per hop\\C 1000 40 60 - B 1000 40 60\\20 20 255 gi:f:1 bi:f:1 cg:Z:20= ...};
\node[font=\scriptsize, anchor=west] at (15.0,1.75) {local GFA (re-aligned)};
\node[draw, circle, inner sep=1.6pt, fill=white] (g1) at (15.3,1.1) {};
\node[draw, circle, inner sep=1.6pt, fill=white] (g2) at (16.0,1.1) {};
\node[draw, circle, inner sep=1.6pt, fill=white] (g3) at (16.7,1.1) {};
\node[draw, circle, inner sep=1.6pt, fill=white] (g4) at (17.4,1.1) {};
\node[draw, circle, inner sep=1.6pt, fill=white] (g5) at (16.35,0.5) {};
\draw (g1) -- (g2) -- (g3) -- (g4);
\draw (g2) -- (g5) -- (g3);
\end{tikzpicture}%
}
\caption{The implicit pangenome graph in impg.
(a) Three sequences A, B and C, which stand for PanSN names such as \texttt{HG1\#1\#chr1}, and two pairwise alignments.
Alignment $a_1$ lies on the plus strand and has a 10 bp insertion in A, and $a_2$ lies on the minus strand (crossed band).
The alignments imply the variation graph $G(A)$ that seqwish would induce, and impg never builds it.
(b) The index keeps one interval tree per sequence.
Each alignment enters the tree of its target and, with the REVERSED flag $R$, the tree of its query.
An entry stores coordinates, strand, the flag and an offset or record number in the alignment file, and never a CIGAR.
The IMPGIDX2 file stores the names, one block per tree and a trailing tree map, and a query decodes only the trees it touches.
(c) A query at depth $m=2$.
The interval A:40-70 projects through $a_1$ to B:40-60, because the 10 inserted bases of A (light red) have no partner in B.
B:40-60 projects through $a_2$ to C:40-60 on the minus strand, and back to A:40-70, which is already visited.
White arrows show the strand.
(d) Outputs of \texttt{query} for this example.
BED rows are merged, and FASTA holds the input interval and the projected intervals.
PAF has one row per hop with its CIGAR slice, after touching or overlapping rows of one sequence pair merge.
The local GFA comes from aligning the intervals again.
File rows in b show query, start, end, strand, target, start, end and CIGAR.
Record $a_1$ is the fixture at tests/test\_transitive\_integrity.rs:462, and $a_2$ is not a fixture.
Coordinates in c and d follow the code with \texttt{-x -m 2 -d 0} and \texttt{-{}-min-transitive-len 0} (derived, not run).}
\label{fig:overview}
\end{figure*}
% src (legend): tests/test_transitive_integrity.rs:462 (a1 fixture); src/impg.rs:142-162 (inverted CIGAR), 164-211 (entry fields), 1686-1712 (two entries), 1758-1954 (IMPGIDX2 and lazy load), 2379-2628 (BFS), 2791-2929 (projection); src/paf.rs:171 (PAF CIGAR offset), src/tpa_parser.rs:116 (TPA record number), src/onealn.rs:697 (1ALN alignment index); src/main.rs:12055-12098 (BED), 12204 and 12648-12760 (PAF merge of contiguous or overlapping rows), 12236-12295 (PAF fields, gi:f and bi:f before cg:Z), 7629 and 7641-7643 (input interval removed for BEDPE and PAF only), 7749-7760 and 12487-12500 (FASTA keeps the input interval), 12324-12340 (GFA by re-alignment); impg query --help (-x, -d required, -m default 2); src/main.rs:4259; hop coordinates derived by hand from src/impg.rs:2791-2929

\textbf{Index.}
impg keeps one interval tree from the coitrees library \citep{Jones2024coitrees} per target sequence (Fig. \ref{fig:overview}b).
% src: src/impg.rs:226, 498-508, 1729-1737; Cargo.lock (coitrees 0.4.0)
An entry holds the other sequence, both coordinate pairs as 32-bit integers, a file number, a data length and one 64-bit word.
% src: src/impg.rs:164-174
Bit 63 of the word is the strand, bit 62 a REVERSED flag, and the rest the CIGAR byte offset or the tracepoint record number.
% src: src/impg.rs:178-207; src/alignment_record.rs:3-10
For bgzipped PAF the offset is a BGZF virtual offset.
% src: src/paf.rs:235-253, 286-298
The index stores no CIGAR and no sequence.
% src: src/impg.rs:164-174
A record with different query and target enters the tree of its target and, flagged REVERSED with swapped coordinates, the tree of its query.
% src: src/impg.rs:1686-1712
Self-alignments enter once, and \texttt{-{}-unidirectional} turns the second insertion off.
% src: src/impg.rs:1688; impg index --help (--unidirectional)
A query on any sequence therefore reaches every alignment that touches it.
% src: src/impg.rs:1686-1712; tests/test_transitive_integrity.rs:226-290
The index refers to alignment files by number, so a query must list them in the order used to build it.
% src: src/impg.rs:1657-1675 (file_index from enumerate); plan.md section 9 (code fix C4)
This paper uses one combined index (\texttt{-i}), not the per-file mode that 100 or more files select by default.
% src: src/main.rs:11277-11291; plan.md section 10

\textbf{On-disk format.}
An IMPGIDX2 file holds an 8-byte magic, the offset of a trailing tree map, the sequence names and lengths, and one block per tree.
% src: src/impg.rs:1758-1825
Loading reads the names and the tree map, and a query decodes a tree when it first touches its sequence.
% src: src/impg.rs:1827-1888, 1891-1954
Decoded trees stay in memory until the process ends, except in \texttt{stats}.
% src: src/impg.rs:1861-1865; eviction only at src/main.rs:13302, 13374

\textbf{Projection.}
For PAF input, impg reads the CIGAR at the stored offset and walks it once, up to the end of the requested range.
% src: src/impg.rs:599-656, 2829-2835
A REVERSED entry walks the inverted CIGAR, with \texttt{I} and \texttt{D} swapped and the order reversed on the minus strand.
% src: src/impg.rs:142-162, 651-654
TPA \citep{Guarracino2026tpa} and 1ALN \citep{Myers2025fastga,Durbin2024onecode} files store tracepoints \citep{Kaushan2026tracepoints}.
% src: citation-keys.md (Guarracino2026tpa, Myers2025fastga, Durbin2024onecode, Kaushan2026tracepoints); src/tpa_parser.rs; src/onealn.rs
For them impg finds the overlapping tracepoint segments, fetches their sequence from FASTA or AGC \citep{Deorowicz2023agc} and rebuilds their CIGAR with WFA2 \citep{MarcoSola2021wfa}.
% src: src/impg.rs:43-71 (thread-local WFA2), 750-927 (segment scan), 1031-1204 (subset rebuild), 1206-1362; src/sequence_index.rs:15-17, 78-90
A hit whose tracepoints cannot be fetched or projected is dropped with a debug message only.
% src: src/impg.rs:1230-1236, 1247-1254, 1330-1338
\pending{E1}{percentage of identical sorted rows between TPA exact and PAF projection of the same records at -m 0, 1 and 2, and the number of dropped tracepoint hits}
\texttt{-{}-approximate} interpolates the boundary inside the first and last overlapping segment, without sequences or WFA2, for BED and BEDPE output.
% src: src/impg.rs:1421-1637; src/main.rs:7393-7408; impg query --help (--approximate)
\pending{E1}{boundary difference in bp between approximate and exact mode, and the wall time of both}
Algorithm \ref{alg:bfs} projects each depth in parallel threads and applies the results one by one in frontier order.
% src: src/impg.rs:2452-2533 (rayon par_iter with ordered collect), 2539-2591 (sequential loop)
\pending{E2}{whether query output is identical at -t 1, 4, 16 and 48}

\textbf{Outputs.}
\texttt{query} writes BED, BEDPE, PAF, FASTA, GFA, VCF and MAF, and PAF rows keep the CIGAR slice of each hop (Fig. \ref{fig:overview}d).
% src: impg query --help (Output formats); src/main.rs:7614-7760, 12276-12295 (PAF fields cg:Z, gi:f, bi:f, an:Z)
Only BED output and one POA-engine GFA case have tests on the alignment backend.
% src: VERIFY_CORRECTIONS.md (core-index-projection, query output formats); tests/test_transitive_integrity.rs; tests/test_graph_output_crush.rs:37-124 (only alignment-backend GFA output test, POA engine); tests/test_aligner_validation.rs:3-31 (-o gfa only for option rejection)
With \texttt{-o gfa}, impg merges the projected intervals, drops their CIGARs and builds a graph of their sequences from scratch.
% src: src/main.rs:12324-12340; src/lib.rs:926-951
The default engine aligns the intervals with each other, induces a graph with an embedded seqwish \citep{Garrison2023seqwish}, smooths it \citep{Garrison2025smoothxg}, normalizes it with GFAffix \citep{Doerr2026gfaffix} and sorts it \citep{Garrison2026gfasort}.
% src: src/lib.rs:44-45, 1112-1131; src/smooth.rs:1-5 (smoothxg-style smoothing); impg query --help (--gfa-engine default pggb); citation-keys.md:17, 19
These graphs come from new alignments, so they need not match any part of $G(A)$.
% src: src/main.rs:12324-12340; plan.md section 11
The embedded seqwish also drops matches shorter than 23 bp by default, which violates the condition of Proposition \ref{prop:restrict}.
% src: impg query --help (--min-match-len default 23)
\authornote{Decision D4 (code fix C15): build -o gfa graphs from the projected PAF records, so that Proposition 2 can apply to them, or keep re-alignment and say so.}

\textbf{Partition and refine.}
\texttt{partition} picks unassigned bases, by default the longest unassigned region, and cuts them into windows of \texttt{-w} bp.
% src: src/commands/partition.rs:754-797 (longest mode), 948-972 (windows); impg partition --help (--selection-mode default longest)
A tail shorter than \texttt{-w} joins the previous window.
% src: src/commands/partition.rs:957-966
It computes the closure of each window with assigned bases masked.
% src: src/commands/partition.rs:369-397 (masked BFS)
Each closure is merged at \texttt{-d} and extended to sequence ends closer than 3,000 bp, and its unmasked part becomes one partition.
% src: src/commands/partition.rs:421-444 (merge_overlaps, extend_to_close_boundaries), 446-451 (mask), 1428-1466; impg partition --help (-d required, --min-boundary-distance default 3000)
By default, singleton partitions are then moved into a flanking partition.
% src: src/commands/partition.rs:31-46, 673-678; impg partition --help (--no-rehome-singletons)
The loop ends for any window of at least 1 bp (Supplementary Note S4).
% src: src/commands/partition.rs:1037-1380; Supplementary Note S4
\pending{E9}{partitions, uncovered bp and bp assigned to more than one partition on ScRAP 142}
\texttt{refine} widens a locus in steps of 1,000 bp, by default up to half its length per side.
% src: impg refine --help (--extension-step 1000, --max-extension 0.5)
It keeps the extension that maximizes the sequences, haplotypes or samples (\texttt{-{}-pansn-mode}) that span both ends by at least 1,000 bp.
% src: impg refine --help (--span-bp 1000, --pansn-mode); src/main.rs:7933-7935 (default Sequence); src/commands/refine.rs:188-273
\pending{E10}{haplotypes spanning both ends before and after refine at TERT and IGKV6-21, and the bp added per side}

\textbf{Memory.}
A build holds every alignment record (56 bytes of fields) and every tree entry (48 bytes) at once.
% src: src/alignment_record.rs:12-21 (2 x u32 + 6 x 8-byte fields = 56 bytes); src/impg.rs:164-174 (40 bytes) plus coitrees Interval first and last (8 bytes); src/impg.rs:1657-1737
For the 389,234,668 HPRC release 2 records, if none is a self-alignment, that is at least 59,163,669,536 bytes by field sizes, not by measurement.
% src: plan.md section 4 (records summed from the 466 TPA headers); 389,234,668 x 56 + 778,469,336 x 48 = 59,163,669,536, recomputed with python3 in this step
\pending{E2}{index build wall time, CPU time, peak RSS and index bytes for the 466 TPA files on the named node}
A code comment puts one AGC decompressor slot at 3.70 GB for the 3.30 GB HPRC archive, with up to 8 slots by default.
% src: src/agc_index.rs:58-83 (default 8, capped by the thread count), 87-100 (comment, 3.70 GB per slot)
Coordinates are 32-bit and operation lengths 29-bit, so impg does not support sequences above 2,147,483,647 bp or operations above 536,870,911 bp.
% src: src/impg.rs:81-109 (29-bit op length), 164-174 (i32); plan.md section 9 (code fix C13, no check today)
